% Part: normal-modal-logic
% Chapter: syntax-and-semantics
% Section: substitution

\documentclass[../../../include/open-logic-section]{subfiles}

\begin{document}

\olfileid{nml}{syn}{sub}

\olsection{যুগপৎ প্রতিস্থাপন}

!!a{formula}~$!A$-এর !!a{propositional variable}-এর সব উপস্থিতি
অন্য কোনো !!{formula} দিয়ে বদলে যে ফল পাওয়া যায়, তা $!A$-এর
একটি \emph{উদাহরণ}। বৈধতা ও !!{derivability} আলোচনায় !!{formula}s-এর
উদাহরণ বারবার আসবে। তাই ধারণাটি নির্ভুলভাবে সংজ্ঞায়িত করা দরকার।

\begin{defn}\ollabel{def:subst-inst}
  ধরা যাক, $!A$ একটি মোডাল !!{formula}, যার সব !!{propositional
    variable}s $p_1$, \dots, $p_n$-এর মধ্যে রয়েছে, এবং $!D_1$,
  \dots, $!D_n$-ও মোডাল !!{formula}s। তাহলে $!A$-তে প্রতিটি
  $p_i$-এর স্থানে যুগপৎ $!D_i$ প্রতিস্থাপন করে পাওয়া ফলকে
  $\SSubst{!A}{\subst{!D_1}{p_1}, \dots, \subst{!D_n}{p_n}}$ দিয়ে
  সংজ্ঞায়িত করি। আনুষ্ঠানিকভাবে এটি $!A$-এর ওপর আবেশের সংজ্ঞা:
  \begin{enumerate}
    \tagitem{prvFalse}{\indcase{!A}{\lfalse}{%
        $\SSubst{\indfrm}{\subst{!D_1}{p_1}, \dots,
          \subst{!D_n}{p_n}}$ হলো $\lfalse$।}}{}
    \tagitem{prvTrue}{\indcase{!A}{\ltrue}{%
        $\SSubst{\indfrm}{\subst{!D_1}{p_1}, \dots,
          \subst{!D_n}{p_n}}$ হলো $\ltrue$।}}{}
    \item \indcase{!A}{q}{$q \not\ident p_i$ ($i
      = 1$, \dots,~$n$) হলে $\SSubst{\indfrm}{\subst{!D_1}{p_1}, \dots,
        \subst{!D_n}{p_n}}$ হলো $q$।}
    \item \indcase{!A}{p_i}{$\SSubst{\indfrm}{\subst{!D_1}{p_1},
        \dots, \subst{!D_n}{p_n}}$ হলো $!D_i$।}
    \tagitem{prvNot}{\indcase{!A}{\lnot !B}{%
        $\SSubst{\indfrm}{\subst{!D_1}{p_1}, \dots, \subst{!D_n}{p_n}}$ হলো
        $\lnot \SSubst{!B}{\subst{!D_1}{p_1}, \dots, \subst{!D_n}{p_n}}$।}}{}
    \tagitem{prvAnd}{\indcase{!A}{(!B \land
        !C)}{$\SSubst{\indfrm}{\subst{!D_1}{p_1}, \dots,
          \subst{!D_n}{p_n}}$ হলো
        \[(\SSubst{!B}{\subst{!D_1}{p_1}, \dots, \subst{!D_n}{p_n}} \land
          \SSubst{!C}{\subst{!D_1}{p_1}, \dots, \subst{!D_n}{p_n}}).\]}}{}
    \tagitem{prvOr}{\indcase{!A}{(!B \lor
          !C)}{$\SSubst{\indfrm}{\subst{!D_1}{p_1}, \dots,
          \subst{!D_n}{p_n}}$ হলো
        \[(\SSubst{!B}{\subst{!D_1}{p_1}, \dots, \subst{!D_n}{p_n}} \lor
          \SSubst{!C}{\subst{!D_1}{p_1}, \dots, \subst{!D_n}{p_n}}).\]}}{}
    \tagitem{prvIf}{\indcase{!A}{(!B \lif
          !C)}{$\SSubst{\indfrm}{\subst{!D_1}{p_1}, \dots,
          \subst{!D_n}{p_n}}$ হলো
        \[(\SSubst{!B}{\subst{!D_1}{p_1}, \dots, \subst{!D_n}{p_n}} \lif
          \SSubst{!C}{\subst{!D_1}{p_1}, \dots, \subst{!D_n}{p_n}}).\]}}{}
    \tagitem{prvIff}{\indcase{!A}{(!B \liff
          !C)}{$\SSubst{\indfrm}{\subst{!D_1}{p_1}, \dots,
          \subst{!D_n}{p_n}}$ হলো
        \[(\SSubst{!B}{\subst{!D_1}{p_1}, \dots, \subst{!D_n}{p_n}} \liff
          \SSubst{!C}{\subst{!D_1}{p_1}, \dots, \subst{!D_n}{p_n}}).\]}}{}
    % BN-SRC-333 TARGET-CORRECTION-BEGIN
    \par\smallskip\noindent\textnormal{\small\textbf{উৎস-সংশোধন (BN-SRC-333):} হিমায়িত উৎসে দ্বিশর্তবাচক প্রতিস্থাপনের ক্ষেত্রটি শর্তবাচক সংযোজকের ট্যাগে বাঁধা; এখানে দ্বিশর্তবাচকের ট্যাগ ব্যবহার করা হয়েছে।} % BN-SRC-333 TARGET-CORRECTION-END
    \tagitem{prvBox}{\indcase{!A}{\Box !B}{%
        $\SSubst{\indfrm}{\subst{!D_1}{p_1}, \dots, \subst{!D_n}{p_n}}$ হলো
        $\Box
        \SSubst{!B}{\subst{!D_1}{p_1}, \dots, \subst{!D_n}{p_n}}.$}}{}
    % BN-SRC-334 TARGET-CORRECTION-BEGIN
    \par\smallskip\noindent\textnormal{\small\textbf{উৎস-সংশোধন (BN-SRC-334):} হিমায়িত উৎসে বক্সের ক্ষেত্রটি ভাষার ট্যাগ ছাড়াই দেওয়া; এখানে সেই ট্যাগ যোগ করা হয়েছে।} % BN-SRC-334 TARGET-CORRECTION-END
    \tagitem{prvDiamond}{\indcase{!A}{\Diamond !B}{%
        $\SSubst{\indfrm}{\subst{!D_1}{p_1}, \dots, \subst{!D_n}{p_n}}$ হলো
        $\Diamond
        \SSubst{!B}{\subst{!D_1}{p_1}, \dots, \subst{!D_n}{p_n}}.$}}{}
  \end{enumerate}
  !!{formula} $\SSubst{!A}{\subst{!D_1}{p_1}, \dots,
    \subst{!D_n}{p_n}}$-কে $!A$-এর একটি \emph{প্রতিস্থাপনজাত
    উদাহরণ} বলে।
\end{defn}

\begin{ex}
  ধরা যাক, $!A$ হলো $p_1 \lif \Box(p_1 \land p_2)$, $!D_1$ হলো
  $\Diamond(p_2 \lif p_3)$ এবং $!D_2$ হলো $\lnot\Box p_1$। তাহলে
  $\SSubst{!A}{\subst{!D_1}{p_1}, \subst{!D_2}{p_2}}$ হলো
  \begin{align*}
    \Diamond(p_2 \lif p_3) &
    \lif \Box(\Diamond(p_2 \lif p_3) \land \lnot\Box p_1)
    \intertext{অন্যদিকে $\SSubst{!A}{\subst{!D_2}{p_1}, \subst{!D_1}{p_2}}$ হলো}
    \lnot\Box p_1 &
    \lif \Box(\lnot\Box p_1 \land \Diamond(p_2 \lif p_3))
    \intertext{লক্ষ্য করুন, যুগপৎ প্রতিস্থাপন সাধারণত পর্যায়ক্রমিক প্রতিস্থাপনের
      সমান নয়। যেমন, উপরের $\SSubst{!A}{\subst{!D_1}{p_1}, \subst{!D_2}{p_2}}$-এর
      সঙ্গে $\Subst{(\Subst{!A}{!D_1}{p_1})}{!D_2}{p_2}$ তুলনা করুন, যা হলো:}
    \Diamond(p_2 \lif p_3) & \Subst{\lif \Box(\Diamond(p_2 \lif p_3) \land p_2)}{\lnot\Box p_1}{p_2}, \text{ অর্থাৎ,}\\
    \Diamond(\lnot\Box p_1 \lif p_3) &
    \lif \Box(\Diamond(\lnot\Box p_1
    \lif p_3) \land \lnot\Box p_1)\\
    \intertext{আর $\Subst{(\Subst{!A}{!D_2}{p_2})}{!D_1}{p_1}$-এর সঙ্গে তুলনা করুন:}
    p_1 & \lif \Subst{\Box(p_1 \land \lnot\Box p_1)}{\Diamond(p_2 \lif p_3)}{p_1}, \text{ অর্থাৎ,}\\
    \Diamond(p_2 \lif p_3) &
    \lif \Box(\Diamond(p_2
    \lif p_3) \land \lnot\Box\Diamond(p_2 \lif p_3)).
  \end{align*}
\end{ex}

\end{document}
